% TOP_PROBLEM: {"id": 2613, "title": "Kourovka Notebook Problem 21.104", "category_id": 17, "category": {"id": 17, "name": "group_theory", "display_name": "Group Theory", "description": "Problems about groups, group actions, representations, and related algebraic structures.", "slug": "group-theory", "order_index": 17, "created_at": "2026-07-31T15:26:25.671Z"}, "status": "open", "problem_number": "KOU-21.104", "set_id": 10}
% FIRST_POSTED: 2026-10-01
% ENTRY_KIND: statement_only
% UnsolvedMath ID 2613; source catalogue code KOU-21.104
% !TeX program = lualatex
% Definition and sources only; this file is not an LLM solution attempt.
\documentclass[11pt,a4paper]{article}
\usepackage[margin=25mm]{geometry}
\usepackage{fontspec}
\setmainfont{lmroman10-regular.otf}[BoldFont=lmroman10-bold.otf,
 ItalicFont=lmroman10-italic.otf,BoldItalicFont=lmroman10-bolditalic.otf]
\usepackage{amsmath,amssymb,mathtools}
\usepackage{unicode-math}
\setmathfont{latinmodern-math.otf}
\AtBeginDocument{\let\setminus\smallsetminus}
\usepackage{xurl,hyperref}
\hypersetup{colorlinks=true,urlcolor=blue}
\setcounter{secnumdepth}{0}
\setlength{\parindent}{0pt}
\setlength{\parskip}{5pt}
\setlength{\emergencystretch}{3em}
\begin{document}
\section*{Kourovka Notebook Problem 21.104}\label{2613}
{\small UnsolvedMath ID 2613 · KOU-21.104 · Group Theory · Kourovka Notebook - New Problems, Issue 21}
\subsection{Definitions and mathematical statement}
Let $T$ be an infinite rooted tree in which vertices at level $j$ have the same finite number $d_j\ge2$ of children. A group $G\le\operatorname{Aut}(T)$ acts faithfully by root-preserving automorphisms. For a vertex $v$, its rigid stabilizer $\operatorname{rist}_G(v)$ consists of elements acting trivially outside the descendant subtree rooted at $v$. Let $\operatorname{rist}_G(j)$ be the subgroup generated by the rigid stabilizers of all vertices at level $j$. The group is a branch group if it is transitive on each level and $[G:\operatorname{rist}_G(j)]<\infty$ for every $j$.

Fix a nontrivial word $w\in F_r$ in $r\ge1$ variables, where nontrivial means nonidentity in the free group. Define word iterates by
\[
 e_0(x_1,\ldots,x_r)=x_1,\qquad
 e_{k+1}(x_1,\ldots,x_r)
 =w(e_k(x_1,\ldots,x_r),x_2,\ldots,x_r).
\]
The group satisfies this Engel-type iterated identity if, for each tuple $(g_1,\ldots,g_r)\in G^r$, some integer $m\ge1$ has $e_m(g_1,\ldots,g_r)=1$. The iteration count may depend on the tuple; no common bound is assumed.

The conjecture states that a finitely generated branch group satisfying such an iterated identity must be torsion. Thus each $g\in G$ must satisfy $g^n=1$ for some positive integer $n$ depending on $g$. The assertion quantifies over every nontrivial word, not just the usual commutator word, and asks for finite element orders rather than a uniform group exponent.
\subsection{Short English statement}
Must a finitely generated branch group satisfying any nontrivial Engel-type iterated word identity have every element of finite order?
\subsection{Sources}
\begin{itemize}
\item[S1] ULAM.AI and the UnsolvedMath Contributors, supplied problems.json, record 2613 (KOU-21.104), Kourovka Notebook - New Problems, Issue 21. Catalogue identity and source transcription; CC BY 4.0.
\url{https://huggingface.co/datasets/ulamai/UnsolvedMath}.
\item[S2] The Kourovka Notebook, 21st edition, new problems, Problem 21.104. Expanded formulation of the supplied catalogue statement.
\url{https://arxiv.org/pdf/1401.0300v47}.
\end{itemize}
\end{document}
